\section{An ample difference and eventual sharpness}
\label{sec:ample-difference}

This section proves the exact rank in an additional range by joining two curves of defect one. The construction uses an ordinary-index-zero class on the same surface and does not require nonvanishing of the restricted cup-product determinant considered below.

\begin{proof}[Proof of Theorem~\ref{thm:ample-difference}]
Work over an algebraic closure of the ground field and put
\[
                 H_e=M^e\otimes N^{-1}.
\]
Choose bases of the two rank-one geometric $\mathbf Z_2$ Tate modules. The homomorphisms induced by $\lambda_M$ and $\lambda_N$ are multiplication by scalars $c_M$ and $c_N$, with $2\mid c_M$ and $c_N$ a unit. Additivity of polarization homomorphisms gives the scalar
\[
                         e c_M-c_N
\]
for $\lambda_{H_e}$. It is a unit for every integer $e$. Thus $a(H_e)=0$ whenever $H_e$ is ample.

The ordinary-index-zero case of Theorem~\ref{thm:p10exact}, applied in degree one to $N$ and $H_e$, gives nonempty opens of smooth curves of defect one in $|N^3|$ and $|H_e^3|$. Choose such a curve $C_0$ in the first system. Since $H_e^3$ is very ample, the curves transverse to $C_0$ form a nonempty open in the second system. This open meets its sharp smooth open, so we can choose $C_1$ satisfying both conditions. The curves meet at
\[
                      q=9(N\cdot H_e)>0
\]
distinct points. Their union $C_*=C_0\cup C_1$ is a connected nodal divisor in $|M^{3e}|$.

The normalization sequence and its degree-one cohomology are
\[
 0\longrightarrow\OO_{C_*}\longrightarrow
 \OO_{C_0}\oplus\OO_{C_1}\longrightarrow k^q\longrightarrow0,
\]
\[
 0\longrightarrow k^{q-1}\longrightarrow H^1(C_*,\OO_{C_*})
 \longrightarrow H^1(C_0,\OO_{C_0})\oplus H^1(C_1,\OO_{C_1})
 \longrightarrow0.
\]
The image of $k^2\to k^q$ is the diagonal line. Both sequences commute with Frobenius, which is bijective on $k^{q-1}$. Exactness of the stable part for finite-dimensional semilinear spaces, proved in Appendix~\ref{app:cohomology}, therefore gives
\[
\begin{split}
 f(C_*)&=f(C_0)+f(C_1)+q-1,\\
 p_a(C_*)&=g(C_0)+g(C_1)+q-1.
\end{split}
\]
Here $f(C_*)$ denotes the stable coherent Frobenius dimension on the nodal curve. It follows that $p_a(C_*)-f(C_*)=2$.

The universal divisor over the complete system $|M^{3e}|$ is flat and projective. Near $C_*$ its fibres have $h^0(\OO)=1$: constants give the lower bound, and upper semicontinuity gives the upper bound. Flatness fixes the Euler characteristic, so $h^1(\OO)$ is constant there. Cohomology and base change give a vector bundle $\mathcal H$ of rank $g=p_a(C_*)$, with its relative Frobenius map $F_S^*\mathcal H\to\mathcal H$.

On every geometric fibre the stable rank is the rank of the $g$th iterate. Indeed the ranks of successive images decrease, and equality of two successive ranks makes the image Frobenius-stable over the perfect residue field. The condition that this iterate have rank at least $g-2$ is open, expressed by its minors. It holds at $C_*$ and hence on a nonempty open of the irreducible complete system. The smooth locus is another nonempty open, because $M^{3e}$ is very ample. Their intersection contains smooth curves of defect at most two.

The upper bound in Theorem~\ref{thm:p10bounds} gives defect at least two for the generic smooth member. Thus its defect is exactly two. Finally
\[
 g=9\chi(M)e^2+1,\qquad r_e=f-1=9\chi(M)e^2-2,
\]
as required. The equality is geometric, so it gives the stated generic rank over every perfect ground field.
\end{proof}

\begin{lemma}\label{lem:ample-positive-square}
Let $M$ be ample on an abelian surface. A line bundle $R$ with $R^2>0$ and $M\cdot R>0$ is ample. If $N$ is also ample, then
\[
                         (M\cdot N)^2\ge M^2N^2.
\]
\end{lemma}
\begin{proof}
The abelian Riemann--Roch identities, valid for every line bundle, are
\[
 \chi(R)=R^2/2,\qquad \deg\lambda_R=\chi(R)^2;
\]
the degree is zero when the polarization homomorphism is not finite \citep[Theorem~9.11 and proof, pp.~131--133]{EGM}. Since $M\cdot R>0$, the line bundle $R^{-1}$ has no nonzero section: a nonzero effective divisor has positive intersection with an ample class. The canonical bundle is trivial, so Serre duality gives $h^2(R)=0$. Therefore
\[
                     h^0(R)=\chi(R)+h^1(R)>0.
\]
The same Riemann--Roch identities make $\ker\lambda_R$ finite. An effective line bundle with finite polarization kernel is ample \citep[Proposition~2.22 and proof, p.~25; Lemma~2.21, pp.~24--25]{EGM}. This proves the first assertion.

For the second, set $x=M^2$, $y=M\cdot N$, $z=N^2$, and use additive divisor notation for $R=xN-yM$. Then $M\cdot R=0$ and $R^2=x(xz-y^2)$. If $xz-y^2>0$, Riemann--Roch and Serre duality imply
\[
                   h^0(R)+h^0(R^{-1})\ge\chi(R)>0.
\]
One of $R,R^{-1}$ would have a nonzero effective divisor of intersection zero with $M$, a contradiction. Hence $y^2\ge xz$.
\end{proof}

\begin{proof}[Proof of the numerical bound in Theorem~\ref{thm:ample-difference}]
Write $x=M^2=2\chi(M)$, $y=M\cdot N$, and $z=N^2$. These are positive and $y^2\ge xz$. If
\[
                     e>\frac{y+\sqrt{y^2-xz}}{x},
\]
the class $eM-N$ has positive square $e^2x-2ey+z$ and positive intersection $ex-y$ with $M$. Lemma~\ref{lem:ample-positive-square} makes it ample. Since $\sqrt{y^2-xz}<y$, the stronger rational inequality $e>2y/x=(M\cdot N)/\chi(M)$ also suffices. This proves an explicit eventual range.
\end{proof}

The hypothesis on $N$ is essential to this argument. It is not asserted that every mixed surface admits an ordinary-index-zero ample class. When such a class exists, Theorem~\ref{thm:ample-difference} settles all sufficiently large degrees for every positive-index polarization, including degrees outside the ordered-index ranges. It says nothing about the value of the first-variation determinant at a prescribed sharp source. The unresolved cases include the remaining small degrees and surfaces on which the ordinary Tate homomorphism of every line bundle is divisible by two.
